\section{AION Flow Security Remediation, Load Qualification and Release Gates}
\label{sec:aion-flow-security-load-qualification}

This section records the production-hardening work completed after the formal AION Flow threat
review. It is designed to be appended to the existing AION technical document. Internal engineering
qualification is complete; independent enterprise certification and representative field validation
remain explicitly open.

\subsection{Canonical request authority}

Every preparation, approval, execution, run inspection and control operation, evaluation decision,
template lifecycle change and production operation now requires a short-lived signed desktop
session. The signature binds the HTTP method, exact target including its query, request-body hash,
device, person, workspace, issue time and one-use nonce. Expired, modified and replayed requests fail
closed. The mother brain then resolves that person against the canonical workspace and Organisation
Authority stores. A role, person or workspace asserted by browser code cannot grant permission.

Execution policy is no longer accepted from the canvas. Each workspace has a signed, versioned,
persisted policy with optimistic revision checks and an exact policy receipt. The default permits
local dry-run work only, permits no external writes and has a zero cloud-spend ceiling. A policy
change therefore remains attributable, reversible and separate from a workflow definition.

\subsection{Customer keystore and production dependencies}

Production runtime no longer contains a universal fallback encryption key. It resolves a
customer-provided secret, an owner-only configured key file or a platform keyring entry, and refuses
to start the affected secure path when none exists. Development installations generate their own
owner-only machine secret instead of sharing a repository constant.

Release engineering now maintains a minimal hash-locked Python production manifest. It produces a
CycloneDX software bill of materials and an Ed25519 signature over the lock hash and bill of
materials. The locked set reported no known vulnerabilities during the recorded audit. Legacy ECDSA
and WeasyPrint paths are excluded from this production environment rather than being inherited from
the all-purpose research installation.

\subsection{Sustained-load evidence}

The earlier timing check measured only repeated graph hashing and was not sufficient. The replacement
qualification covers three distinct synthetic profiles. A 1,500-node linear graph with 1,499 edges
passed production shape, identity and authority validation with a recorded 95th-percentile time of
997.635 milliseconds. A 256-route fan-out and fan-in graph with 512 edges produced deterministic
ordered results with a 95th-percentile processing time of 38.646 milliseconds. A simulated
long-running workflow wrote 250 atomic local checkpoints, recovered checkpoint 250 and recorded a
95th-percentile checkpoint-write time of 0.625 milliseconds.

The run used synthetic identifiers, made no external model calls and performed no external business
actions. Its signed report hash is
\texttt{f1f1d047cbc628f0}\allowbreak\texttt{2be7715fbc2c8f93}\allowbreak
\texttt{fcde9fff5b18b993}\allowbreak\texttt{03d4977ecd659263}.
These figures qualify the tested machine and build, not every future customer device. The Production
centre now exposes the same signed qualification so each installation can generate its own evidence.

\subsection{Accessibility and privacy boundary}

The Production centre is exposed as a labelled modal dialog with semantic controls, a polite live
status region, keyboard focus on entry, Escape dismissal and focus restoration. The mobile surface
remains review-only. Encrypted imports restore neither credentials nor execution authority.
Operational telemetry contains hashes, counts, node families, status totals and latency buckets; it
does not contain workflow content, credentials or customer identifiers. Load evidence likewise uses
synthetic content only.

The governing privacy rule remains that the canvas, imported graph, model output, evidence and
provider response are untrusted inputs. Credentials remain opaque Visual Vault references. Persisted
customer policy determines disclosure, and only an authorised capability plus a verified receipt can
prove an external effect.

\subsection{Independent release work still required}

The implementation team cannot self-certify the remaining release gates. Before an applicable
enterprise release, AION Flow still requires independent penetration testing; a WCAG~2.2 AA audit
including keyboard, screen-reader, zoom, contrast and cognitive-accessibility testing;
jurisdiction-specific privacy, employment, retention, AI-governance and sector review; dependency
licence and update-signature review; platform signing and notarisation with the release owner's
identities; and representative trials across real organisations, roles, workspaces, private models,
customer-cloud endpoints and disaster-recovery exercises.

Accordingly, remediation tasks \texttt{AFLOW-1107R1} through \texttt{AFLOW-1107R4} and sustained-load
task \texttt{AFLOW-1108} are complete. Task \texttt{AFLOW-1109} remains open only for the independent,
jurisdictional and representative field evidence described above. No model, template or canvas policy
may waive those gates.
